% =========================================================
\section{Übung 1: Zeitbereichanalyse R--L Lasten}
% =========================================================

\subsection*{a) Zeitfunktion des Stromes $i_R(t)$ im stationären Betrieb}

\paragraph{Gegeben}
Periodische Umschaltung zweier Stromschalter $S_1$ und $S_2$ mit
\[
t\in[0,t_\mathrm{ein}] \quad \text{(EIN-Phase)}, 
\qquad
t\in[t_\mathrm{ein},\,t_\mathrm{ein}+t_\mathrm{aus}] \quad \text{(AUS-Phase)}.
\]
R--L Last mit $R$, $L$, Versorgung $U_b$.
(Notation wie im Aufgabenblatt: $i_{R1}$ in EIN-Phase, $i_{R2}$ in AUS-Phase.)

\paragraph{Gesucht}
Zeitfunktion $i_R(t)$ im stationären Betrieb.

\paragraph{Lösungsweg}
\textbf{1) EIN-Phase $t\in[0,t_\mathrm{ein}]$:}
Maschenregel:
\[
R\,i_{R1}(t) + L\frac{di_{R1}}{dt} = U_b.
\]
Homogener Anteil:
\[
R\,i_{R1h}+L\frac{di_{R1h}}{dt}=0
\;\Rightarrow\;
i_{R1h}(t)=C\,e^{-\frac{R}{L}t}.
\]
Partikuläre Lösung (Konstante):
\[
i_{R1p}(t)=\frac{U_b}{R}.
\]
Damit:
\[
i_{R1}(t)=\frac{U_b}{R}+C\,e^{-\frac{R}{L}t}.
\]
Mit Startwert $i_{R1}(0)=i_\mathrm{min}$:
\[
C=i_\mathrm{min}-\frac{U_b}{R}
\quad\Rightarrow\quad
\boxed{
i_{R1}(t)=\frac{U_b}{R}+\left(i_\mathrm{min}-\frac{U_b}{R}\right)e^{-\frac{R}{L}t}
}.
\]

\textbf{2) AUS-Phase $t\in[t_\mathrm{ein},t_\mathrm{ein}+t_\mathrm{aus}]$:}
Maschenregel (Quelle weg):
\[
R\,i_{R2}(t) + L\frac{di_{R2}}{dt} = 0
\;\Rightarrow\;
\boxed{
i_{R2}(t)=D\,e^{-\frac{R}{L}(t-t_\mathrm{ein})}
}.
\]
Mit $i_{R2}(t_\mathrm{ein})=i_{R1}(t_\mathrm{ein})=i_\mathrm{max}$ folgt:
\[
\boxed{
i_{R2}(t)=i_\mathrm{max}\,e^{-\frac{R}{L}(t-t_\mathrm{ein})}
}.
\]

\textbf{3) Stationaritätsbedingung (Periodizität):}
\[
i_{R2}(t_\mathrm{ein}+t_\mathrm{aus}) = i_\mathrm{min}.
\]
Also:
\[
i_\mathrm{min}=i_\mathrm{max}e^{-\frac{R}{L}t_\mathrm{aus}}
\;\Rightarrow\;
\boxed{
i_\mathrm{max}=i_\mathrm{min}\,e^{\frac{R}{L}t_\mathrm{aus}}
}.
\]
Und aus EIN-Phase bei $t=t_\mathrm{ein}$:
\[
i_\mathrm{max}=\frac{U_b}{R}+\left(i_\mathrm{min}-\frac{U_b}{R}\right)e^{-\frac{R}{L}t_\mathrm{ein}}.
\]
Damit sind $i_\mathrm{min},i_\mathrm{max}$ eindeutig bestimmbar und $i_R(t)$ ist stückweise durch $i_{R1},i_{R2}$ gegeben.

% ---------------------------------------------------------

\subsection*{b) Mittelwert $\overline{i_R}$}

\paragraph{Gegeben}
Stückweise Stromverlauf aus Teil a).

\paragraph{Gesucht}
\[
\overline{i_R}=\frac{1}{T}\int_0^T i_R(t)\,dt, \qquad T=t_\mathrm{ein}+t_\mathrm{aus}.
\]

\paragraph{Lösungsweg}
\[
\overline{i_R}=\frac{1}{T}\left(
\int_0^{t_\mathrm{ein}} i_{R1}(t)\,dt + \int_{t_\mathrm{ein}}^{T} i_{R2}(t)\,dt
\right).
\]
Mit
\[
i_{R1}(t)=\frac{U_b}{R}+\left(i_\mathrm{min}-\frac{U_b}{R}\right)e^{-\frac{R}{L}t}
\]
folgt
\[
\int_0^{t_\mathrm{ein}} i_{R1}(t)\,dt
=
\frac{U_b}{R}t_\mathrm{ein}
+
\left(i_\mathrm{min}-\frac{U_b}{R}\right)\frac{L}{R}\left(1-e^{-\frac{R}{L}t_\mathrm{ein}}\right).
\]
Und
\[
\int_{t_\mathrm{ein}}^{T} i_{R2}(t)\,dt
=
\int_0^{t_\mathrm{aus}} i_\mathrm{max}e^{-\frac{R}{L}\tau}\,d\tau
=
i_\mathrm{max}\frac{L}{R}\left(1-e^{-\frac{R}{L}t_\mathrm{aus}}\right).
\]
Damit:
\[
\boxed{
\overline{i_R}=\frac{1}{T}\left[
\frac{U_b}{R}t_\mathrm{ein}
+
\left(i_\mathrm{min}-\frac{U_b}{R}\right)\frac{L}{R}\left(1-e^{-\frac{R}{L}t_\mathrm{ein}}\right)
+
i_\mathrm{max}\frac{L}{R}\left(1-e^{-\frac{R}{L}t_\mathrm{aus}}\right)
\right].
}
\]

% ---------------------------------------------------------

\subsection*{c) Effektivwert $I_{\mathrm{RMS}}$}

\paragraph{Gesucht}
\[
I_{\mathrm{RMS}}=\sqrt{\frac{1}{T}\int_0^T i_R^2(t)\,dt}.
\]

\paragraph{Lösungsweg}
\[
I_{\mathrm{RMS}}^2=\frac{1}{T}\left(
\int_0^{t_\mathrm{ein}} i_{R1}^2(t)\,dt + \int_{t_\mathrm{ein}}^T i_{R2}^2(t)\,dt
\right).
\]
Für die AUS-Phase:
\[
\int_{t_\mathrm{ein}}^T i_{R2}^2(t)\,dt
=
\int_0^{t_\mathrm{aus}} i_\mathrm{max}^2 e^{-2\frac{R}{L}\tau}\,d\tau
=
i_\mathrm{max}^2\frac{L}{2R}\left(1-e^{-2\frac{R}{L}t_\mathrm{aus}}\right).
\]
Für die EIN-Phase setze
\[
i_{R1}(t)=A+B e^{-\frac{R}{L}t},
\quad A=\frac{U_b}{R},\quad B=i_\mathrm{min}-\frac{U_b}{R}.
\]
Dann
\[
i_{R1}^2=A^2+2ABe^{-\frac{R}{L}t}+B^2e^{-2\frac{R}{L}t}.
\]
Also:
\[
\int_0^{t_\mathrm{ein}} i_{R1}^2(t)\,dt
=
A^2 t_\mathrm{ein}
+2AB\frac{L}{R}\left(1-e^{-\frac{R}{L}t_\mathrm{ein}}\right)
+B^2\frac{L}{2R}\left(1-e^{-2\frac{R}{L}t_\mathrm{ein}}\right).
\]
Somit:
\[
\boxed{
I_{\mathrm{RMS}}=\sqrt{\frac{1}{T}\left[
A^2 t_\mathrm{ein}
+2AB\frac{L}{R}\left(1-e^{-\frac{R}{L}t_\mathrm{ein}}\right)
+B^2\frac{L}{2R}\left(1-e^{-2\frac{R}{L}t_\mathrm{ein}}\right)
+i_\mathrm{max}^2\frac{L}{2R}\left(1-e^{-2\frac{R}{L}t_\mathrm{aus}}\right)
\right]}.
}
\]

% ---------------------------------------------------------

\subsection*{d) Spannungen $u_R(t)$ und $u_L(t)$}

\paragraph{Gesucht}
Zeitfunktionen $u_R(t)$, $u_L(t)$.

\paragraph{Lösungsweg}
\[
u_R(t)=R\,i_R(t),
\qquad
u_L(t)=L\frac{di_R}{dt}.
\]
In EIN-Phase aus DGL:
\[
u_L(t)=U_b-u_R(t)=U_b-Ri_{R1}(t).
\]
In AUS-Phase:
\[
u_L(t)=-u_R(t)=-R i_{R2}(t).
\]

% ---------------------------------------------------------

\subsection*{e) Simulation/Verifikation (LTspice/PLECS)}

\paragraph{Lösungsweg}
Parameter wie im Aufgabenblatt setzen, periodische Schalteransteuerung definieren,
Zeitbereich simulieren, stationären Bereich auswerten und $i_\mathrm{min},i_\mathrm{max},\overline{i},I_{\mathrm{RMS}}$
mit den analytischen Ausdrücken aus a)–c) vergleichen.
% =========================================================



% =========================================================
\section{Übung 2: Einpuls-Gleichrichter M1U }
% =========================================================

\subsection*{a) Zeitverlauf $u_2(t)$ und $u_R(t)$ für $t\in[0,60\,\mathrm{ms}]$}

\paragraph{Gegeben}
Einpuls-Gleichrichter M1U, Eingang
\[
u_2(t)=\hat u_2\sin(\omega t),
\quad f_1=50\,\mathrm{Hz}\Rightarrow T=20\,\mathrm{ms}.
\]
Lastspannung
\[
u_R(t)=
\begin{cases}
\hat u_2\sin(\omega t), & 0<\omega t<\pi\\
0, & \pi<\omega t<2\pi
\end{cases}
\quad\text{(periodisch)}.
\]

\paragraph{Gesucht}
Plot über 3 Perioden. (Rechenweg: Stückweise Definition oben.)

\subsection*{b) Mittelwert und Effektivwert von $u_R(t)$}

\paragraph{Gesucht}
\[
\overline{u_R}=\frac{1}{T}\int_0^T u_R(t)\,dt,
\qquad
U_{R,\mathrm{RMS}}=\sqrt{\frac{1}{T}\int_0^T u_R^2(t)\,dt}.
\]

\paragraph{Lösungsweg}
Substitution $\alpha=\omega t$, $dt=\frac{d\alpha}{\omega}$, $T=\frac{2\pi}{\omega}$.

\textbf{Mittelwert:}
\[
\overline{u_R}
=\frac{1}{2\pi}\int_0^{\pi}\hat u_2\sin\alpha\,d\alpha
=\frac{\hat u_2}{2\pi}\left[-\cos\alpha\right]_0^\pi
=\boxed{\frac{\hat u_2}{\pi}}.
\]

\textbf{Effektivwert:}
\[
U_{R,\mathrm{RMS}}
=\sqrt{\frac{1}{2\pi}\int_0^{\pi}\hat u_2^2\sin^2\alpha\,d\alpha}
=
\sqrt{\frac{\hat u_2^2}{2\pi}\cdot\frac{\pi}{2}}
=
\boxed{\frac{\hat u_2}{2}}.
\]

\subsection*{c) Erste sechs Harmonische von $u_R(t)$}

\paragraph{Gesucht}
Fourierreihe
\[
u_R(\alpha)=\frac{a_0}{2}+\sum_{k=1}^{\infty}\left(a_k\cos(k\alpha)+b_k\sin(k\alpha)\right)
\]
mit
\[
a_0=\frac{1}{\pi}\int_0^{2\pi}u_R(\alpha)\,d\alpha,\quad
a_k=\frac{1}{\pi}\int_0^{2\pi}u_R(\alpha)\cos(k\alpha)\,d\alpha,\quad
b_k=\frac{1}{\pi}\int_0^{2\pi}u_R(\alpha)\sin(k\alpha)\,d\alpha.
\]

\paragraph{Lösungsweg}
Da $u_R(\alpha)=\hat u_2\sin\alpha$ nur für $\alpha\in[0,\pi]$, sonst 0:

\[
a_0=\frac{1}{\pi}\int_0^{\pi}\hat u_2\sin\alpha\,d\alpha
=\boxed{\frac{2\hat u_2}{\pi}}.
\]

\[
b_k=\frac{\hat u_2}{\pi}\int_0^{\pi}\sin\alpha\sin(k\alpha)\,d\alpha
=
\begin{cases}
\frac{\hat u_2}{2}, & k=1\\
0, & k\neq 1
\end{cases}
\]

\[
a_k=\frac{\hat u_2}{\pi}\int_0^{\pi}\sin\alpha\cos(k\alpha)\,d\alpha
=
\frac{\hat u_2}{\pi}\cdot\frac{2}{1-k^2}
\quad\text{für gerade }k,
\qquad
a_k=0\ \text{für ungerade }k.
\]

Damit (erste \textbf{sechs} Nicht-Null-Koeffizienten wie in Lösung):
\[
\boxed{
a_2=-\frac{2}{3\pi}\hat u_2,\ 
a_4=-\frac{2}{15\pi}\hat u_2,\ 
a_6=-\frac{2}{35\pi}\hat u_2,\ 
a_8=-\frac{2}{63\pi}\hat u_2,\ 
a_{10}=-\frac{2}{99\pi}\hat u_2,\ 
a_{12}=-\frac{2}{143\pi}\hat u_2,\ 
b_1=\frac{\hat u_2}{2}
}.
\]
% =========================================================



% =========================================================
\section{Übung 3: Zweipuls-Gleichrichter}
% =========================================================

\subsection*{a) Zeitverlauf $u_2(t)$ und $u_R(t)$ über drei Perioden}

\paragraph{Gegeben}
Eingang:
\[
u_2(t)=\hat u_2\sin(\omega t),\quad f_1=50\,\mathrm{Hz}.
\]
Vollwellige Gleichrichtung:
\[
u_R(t)=|u_2(t)|=\hat u_2|\sin(\omega t)|.
\]

\paragraph{Gesucht}
Zeitverlauf für $t\in[0,60\,\mathrm{ms}]$.

\paragraph{Lösungsweg}
Stückweise:
\[
u_R(t)=
\begin{cases}
\hat u_2\sin(\omega t), & 0<\omega t<\pi\\
-\hat u_2\sin(\omega t), & \pi<\omega t<2\pi
\end{cases}
\quad\text{periodisch}.
\]

\subsection*{b) Mittelwert und Effektivwert von $u_R(t)$}

\paragraph{Gesucht}
\[
\overline{u_R},\quad U_{R,\mathrm{RMS}}.
\]

\paragraph{Lösungsweg}
Mit $\alpha=\omega t$:

\textbf{Mittelwert:}
\[
\overline{u_R}
=\frac{1}{2\pi}\int_0^{2\pi}\hat u_2|\sin\alpha|\,d\alpha
=\frac{2}{2\pi}\int_0^{\pi}\hat u_2\sin\alpha\,d\alpha
=\boxed{\frac{2\hat u_2}{\pi}}.
\]

\textbf{Effektivwert:}
\[
U_{R,\mathrm{RMS}}
=\sqrt{\frac{1}{2\pi}\int_0^{2\pi}\hat u_2^2\sin^2\alpha\,d\alpha}
=
\sqrt{\frac{\hat u_2^2}{2\pi}\cdot\pi}
=
\boxed{\frac{\hat u_2}{\sqrt{2}}}.
\]

\subsection*{c) Harmonische (Fourier-Koeffizienten)}

\paragraph{Lösungsweg}
Da $u_R(\alpha)=\hat u_2|\sin\alpha|$ gerade ist, gilt:
\[
b_k=0.
\]
Koeffizienten:
\[
a_0=\frac{1}{\pi}\int_0^{2\pi}\hat u_2|\sin\alpha|\,d\alpha=\boxed{\frac{4\hat u_2}{\pi}},
\]
\[
a_k=\frac{1}{\pi}\int_0^{2\pi}\hat u_2|\sin\alpha|\cos(k\alpha)\,d\alpha
=
\begin{cases}
-\frac{4\hat u_2}{\pi}\frac{1}{k^2-1}, & k\ \text{gerade}\\
0, & k\ \text{ungerade}
\end{cases}
\]
(erste relevante gerade $k=2,4,6,\dots$).
% =========================================================



% =========================================================
\section{Übung 4: Wechselstromsteller W1 (R-Last)}
% =========================================================

\subsection*{a) Zeitverlauf von $u_R(t)$ und $i_R(t)$ für $t\in[0,60\,\mathrm{ms}]$}

\paragraph{Gegeben}
Wechselstromsteller (Phasenanschnitt) mit ohmscher Last $R$.
Eingangsspannung:
\[
u_2(t)=\hat U \sin(\omega t),\qquad \omega=2\pi f.
\]
Zündwinkel $\alpha$ pro Halbwelle.

\paragraph{Gesucht}
Zeitverlauf von Lastspannung $u_R(t)$ und Laststrom $i_R(t)$ über drei Perioden.

\paragraph{Lösungsweg}
Für eine reine R-Last gilt:
\[
i_R(t)=\frac{u_R(t)}{R}.
\]
Phasenanschnitt bedeutet: Leitung nur im Intervall $\omega t\in[\alpha,\pi]$ (positive Halbwelle)
und $\omega t\in[\pi+\alpha,2\pi]$ (negative Halbwelle).

Definiere $\beta=\omega t$ (periodisch mit $2\pi$):
\[
u_R(\beta)=
\begin{cases}
\hat U \sin\beta, & \alpha \le \beta \le \pi,\\[4pt]
\hat U \sin\beta, & \pi+\alpha \le \beta \le 2\pi,\\[4pt]
0, & \text{sonst.}
\end{cases}
\]
und damit:
\[
i_R(\beta)=\frac{u_R(\beta)}{R}.
\]

% ---------------------------------------------------------

\subsection*{b) Mittelwert und Effektivwert von $u_R(t)$ und $i_R(t)$}

\paragraph{Gegeben}
Stückweise Definition aus a), $\beta=\omega t$.

\paragraph{Gesucht}
\[
\overline{u_R},\quad U_{R,\mathrm{RMS}},\quad \overline{i_R},\quad I_{R,\mathrm{RMS}}.
\]

\paragraph{Lösungsweg}
\textbf{Mittelwert:}
Wegen Symmetrie (positive und negative Halbwelle gleich gross, Vorzeichen wechselt):
\[
\boxed{\overline{u_R}=0}
\qquad\Rightarrow\qquad
\boxed{\overline{i_R}=0}.
\]

\textbf{Effektivwert:}
\[
U_{R,\mathrm{RMS}}=\sqrt{\frac{1}{2\pi}\int_0^{2\pi}u_R^2(\beta)\,d\beta}.
\]
Da $u_R(\beta)=\hat U\sin\beta$ nur in den Leitintervallen:
\[
U_{R,\mathrm{RMS}}^2=\frac{1}{2\pi}\left(
\int_{\alpha}^{\pi}\hat U^2\sin^2\beta\,d\beta
+
\int_{\pi+\alpha}^{2\pi}\hat U^2\sin^2\beta\,d\beta
\right)
=\frac{\hat U^2}{\pi}\int_{\alpha}^{\pi}\sin^2\beta\,d\beta.
\]
Mit
\[
\sin^2\beta=\frac{1-\cos(2\beta)}{2}
\]
folgt
\[
\int_{\alpha}^{\pi}\sin^2\beta\,d\beta
=
\left[\frac{\beta}{2}-\frac{\sin(2\beta)}{4}\right]_{\alpha}^{\pi}
=
\frac{\pi-\alpha}{2}+\frac{\sin(2\alpha)}{4}.
\]
Somit:
\[
\boxed{
U_{R,\mathrm{RMS}}
=
\hat U \sqrt{\frac{1}{\pi}\left(\frac{\pi-\alpha}{2}+\frac{\sin(2\alpha)}{4}\right)}
=
\hat U\sqrt{\frac{\pi-\alpha}{2\pi}+\frac{\sin(2\alpha)}{4\pi}}
}.
\]
Für den Strom:
\[
\boxed{I_{R,\mathrm{RMS}}=\frac{U_{R,\mathrm{RMS}}}{R}}.
\]

% ---------------------------------------------------------

\subsection*{c) Wirkleistung $P$ an der ohmschen Last}

\paragraph{Gesucht}
Wirkleistung $P$.

\paragraph{Lösungsweg}
Für reine R-Last:
\[
P = R\, I_{R,\mathrm{RMS}}^2=\frac{U_{R,\mathrm{RMS}}^2}{R}.
\]
Mit obigem Ergebnis:
\[
\boxed{
P
=
\frac{\hat U^2}{R}
\left(
\frac{\pi-\alpha}{2\pi}+\frac{\sin(2\alpha)}{4\pi}
\right)
}.
\]

% ---------------------------------------------------------

\subsection*{d) Fourierkoeffizienten von $i_R(t)$ (Spektrum)}

\paragraph{Gesucht}
Fourierreihe von $i_R(\beta)$ und Struktur der Koeffizienten.

\paragraph{Lösungsweg}
\[
i_R(\beta)=\frac{u_R(\beta)}{R}
\]
ist ein \textbf{ungerades} Signal (anti-symmetrisch), daher gilt:
\[
\boxed{a_k=0\ \text{für alle }k,\qquad a_0=0.}
\]
Es bleiben nur Sinuskoeffizienten:
\[
b_k=\frac{1}{\pi}\int_0^{2\pi} i_R(\beta)\sin(k\beta)\,d\beta
=
\frac{2}{\pi}\int_{\alpha}^{\pi} \frac{\hat U}{R}\sin\beta\sin(k\beta)\,d\beta.
\]
Mit Produkt-zu-Summe:
\[
\sin\beta\sin(k\beta)=\frac{1}{2}\left[\cos((k-1)\beta)-\cos((k+1)\beta)\right].
\]
Dann:
\[
b_k=\frac{\hat U}{\pi R}\int_{\alpha}^{\pi}\left[\cos((k-1)\beta)-\cos((k+1)\beta)\right]\,d\beta.
\]
Integration:
\[
\boxed{
b_k
=
\frac{\hat U}{\pi R}
\left[
\frac{\sin((k-1)\beta)}{k-1}
-
\frac{\sin((k+1)\beta)}{k+1}
\right]_{\alpha}^{\pi}
}.
\]
Damit ist die Fourierreihe:
\[
\boxed{
i_R(\beta)=\sum_{k=1}^{\infty} b_k \sin(k\beta).
}
\]

% ---------------------------------------------------------

\subsection*{e) Blindleistung und Verzerrungsblindleistung}

\paragraph{Gesucht}
Grundschwingungsblindleistung und Verzerrungsblindleistung.

\paragraph{Lösungsweg}
Bei ohmscher Last ist die \textbf{Grundschwingungsblindleistung} ideal:
\[
\boxed{Q_1=0}
\]
(Grundschwingungsstrom und -spannung in Phase).

Nichtsinusförmige Ströme erzeugen jedoch Verzerrungsblindleistung:
\[
S = U_{\mathrm{RMS}} I_{\mathrm{RMS}},\qquad P=R I_{\mathrm{RMS}}^2,
\]
\[
\boxed{D=\sqrt{S^2-P^2-Q_1^2}=\sqrt{S^2-P^2}}.
\]



% =========================================================
\section{Übung 5: Gleichstromsteller (Schalter) }
% =========================================================

\subsection*{a) $i_L(t)$ und $u_R(t)$ über eine Schaltperiode}

\paragraph{Gegeben}
Gleichstromquelle $U_{DC}$, Schaltperiode $T_s=T_e+T_a$.
Einschaltzeit $T_e$, Ausschaltzeit $T_a$.
R--L Last ($R_1$, $L_1$). Zeitkonstante:
\[
\boxed{\tau=\frac{L_1}{R_1}}.
\]

\paragraph{Gesucht}
Stromverlauf $i_L(t)$ und Lastspannung $u_R(t)$ im stationären Betrieb.

\paragraph{Lösungsweg}
\textbf{Intervall 1 (EIN): $0<t<T_e$}
\[
L_1\frac{di_L}{dt}+R_1 i_L = U_{DC}.
\]
Allg. Lösung:
\[
\boxed{
i_L(t)=\frac{U_{DC}}{R_1}+\left(i_{\min}-\frac{U_{DC}}{R_1}\right)e^{-t/\tau}
},\qquad 0<t<T_e.
\]
Endwert:
\[
\boxed{
i_{\max}=i_L(T_e)=\frac{U_{DC}}{R_1}+\left(i_{\min}-\frac{U_{DC}}{R_1}\right)e^{-T_e/\tau}.
}
\]

\textbf{Intervall 2 (AUS): $T_e<t<T_s$}
Quelle weg:
\[
L_1\frac{di_L}{dt}+R_1 i_L = 0.
\]
Damit:
\[
\boxed{
i_L(t)= i_{\max}\,e^{-(t-T_e)/\tau}
},\qquad T_e<t<T_s.
\]
Endwert:
\[
\boxed{
i_{\min}=i_L(T_s)=i_{\max}\,e^{-T_a/\tau}.
}
\]

\textbf{Stationarität (Periodenbedingung):}
\[
i_{\min}=i_{\max}e^{-T_a/\tau}
\quad\Rightarrow\quad
\boxed{T_a=-\tau\ln\!\left(\frac{i_{\min}}{i_{\max}}\right)}.
\]

\textbf{Lastspannung:}
\[
\boxed{u_R(t)=R_1\, i_L(t)}.
\]

% ---------------------------------------------------------

\subsection*{b) Bestimmung einer Schaltzeit aus $i_{\min}$, $i_{\max}$ (typische Prüfungsform)}

\paragraph{Gegeben}
$i_{\min}$ und $i_{\max}$ (aus Diagramm/Simulation).

\paragraph{Gesucht}
$T_a$ (oder $T_e$).

\paragraph{Lösungsweg}
Direkt aus der AUS-Phase:
\[
i_{\min}=i_{\max}e^{-T_a/\tau}
\;\Rightarrow\;
\boxed{T_a=-\tau\ln\!\left(\frac{i_{\min}}{i_{\max}}\right)}.
\]
Dann $T_e=T_s-T_a$ (falls $T_s$ gegeben).


% =========================================================
\section{Übung 6: PWM-Wechselrichter }
% =========================================================

\subsection*{a) Herleitung von $i_L(t)$ und $u_L(t)$ in einem PWM-Intervall}

\paragraph{Gegeben}
PWM-Schaltsignal erzeugt zeitlich wechselnde Lastspannung.
In jedem Intervall gilt entweder:
\[
u_L(t)=U_{DC}\quad \text{oder}\quad u_L(t)=0
\]
(idealisiert, je nach Schalterzustand).
R--L Last ($R$, $L$), $\tau=\frac{L}{R}$.

\paragraph{Gesucht}
Stromverlauf $i_L(t)$ stückweise und Darstellung über mehrere Schaltintervalle.

\paragraph{Lösungsweg}
Betrachte ein beliebiges Intervall $[t_i,t_{i+1}]$.

\textbf{Fall 1: $u_L(t)=U_{DC}$ auf $[t_i,t_{i+1}]$}
\[
L\frac{di}{dt}+Ri=U_{DC}.
\]
Mit Anfangswert $i(t_i)$:
\[
\boxed{
i(t)=\frac{U_{DC}}{R}+\left(i(t_i)-\frac{U_{DC}}{R}\right)e^{-(t-t_i)/\tau}
},\qquad t\in[t_i,t_{i+1}].
\]

\textbf{Fall 2: $u_L(t)=0$ auf $[t_i,t_{i+1}]$}
\[
L\frac{di}{dt}+Ri=0
\Rightarrow
\boxed{
i(t)=i(t_i)e^{-(t-t_i)/\tau}
},\qquad t\in[t_i,t_{i+1}].
\]

Die Endwerte $i(t_{i+1})$ sind jeweils die Anfangswerte des nächsten Intervalls.

\textbf{Lastspannung:}
Je nach PWM-Zustand:
\[
\boxed{u_L(t)\in\{0,\ U_{DC}\}}.
\]

% ---------------------------------------------------------

\subsection*{b) Wodurch wird die Hauptfrequenz des Laststroms bestimmt?}

\paragraph{Gesucht}
Hauptfrequenz-Aussage.

\paragraph{Lösungsweg}
Die PWM enthält zwei charakteristische Frequenzen:
\begin{itemize}
\item Grundfrequenz $f_1$ der Referenz (Sinus)
\item Schaltfrequenz $f_s$ des Trägers
\end{itemize}
Die \textbf{Hauptfrequenz} der Laststromgrundschwingung wird durch die Referenz bestimmt:
\[
\boxed{f_{i,\mathrm{main}}=f_1.}
\]
Die Schaltfrequenz erzeugt Oberschwingungen und Ripple um Vielfache von $f_s$.

% ---------------------------------------------------------

\subsection*{c) Vergleich idealer und realer Laststrom (qualitativ)}

\paragraph{Gesucht}
Warum weicht realer PWM-Strom vom idealen ab?

\paragraph{Lösungsweg}
Abweichungen entstehen durch:
\begin{itemize}
\item endliche Schaltfrequenz $f_s$ (Ripple)
\item nichtideale Schalter (Spannungsabfälle, Totzeit)
\item Filterwirkung der Last ($R$--$L$) dämpft hohe Harmonische
\end{itemize}



% =========================================================
\section{Übung 7: Gesteuerter Gleichrichter + Gleichstrommaschine}
% =========================================================

\subsection*{a) Warum läuft die Maschine unterhalb einer Winkelgrenze nicht sauber an?}

\paragraph{Gegeben}
Gesteuerter Gleichrichter speist fremderregte Gleichstrommaschine.
Induzierte Spannung:
\[
\boxed{E = k\Phi \omega.}
\]
Gleichrichter-Ausgangsspannung $U_a(\alpha)$ ist abhängig vom Steuerwinkel $\alpha$.

\paragraph{Gesucht}
Erklärung der mechanischen Schwingungen und Nicht-Anlauf für kleine $\alpha$.

\paragraph{Lösungsweg}
Ankerstrom fliesst nur, wenn die momentane Gleichrichterspannung die induzierte Spannung übersteigt:
\[
u_a(t) > E.
\]
In der Lösung wird ein Grenzwinkel $\beta$ eingeführt mit:
\[
\boxed{E = U_{2m}\sin\beta}.
\]
Dann gilt:
\[
\boxed{\text{Stromfluss nur für }\alpha>\beta.}
\]
Falls $\alpha<\beta$:
\[
i_a(t)\approx 0\Rightarrow M\approx 0
\]
und der Motor erzeugt kein stabiles Drehmoment, was zu Stillstand bzw. Schwingungen führt.

% ---------------------------------------------------------

\subsection*{b) Warum steigt $n$ mit $\alpha$ zuerst, erreicht ein Maximum und fällt dann wieder?}

\paragraph{Gegeben}
Drehmoment:
\[
\boxed{M=I_a\Phi}
\]
und bei Fremderregung $\Phi=\mathrm{konstant}$.
Damit bestimmt $I_a$ das Drehmoment.

\paragraph{Gesucht}
Erklärung des nichtmonotonen Zusammenhangs $n(\alpha)$.

\paragraph{Lösungsweg}
Mittelwert-Ankerstrom (wie in der Lösung) über das leitende Intervall:
\[
\boxed{
I_{a,\mathrm{avr}}
=
\frac{1}{\pi R_a}\int_{\alpha}^{\pi-\beta}\left(U_{2m}\sin\gamma - E\right)\,d\gamma
}
\quad \text{mit } \gamma=\omega t.
\]
Für kleine $\alpha$ ist das Leitintervall gross, $I_{a,\mathrm{avr}}$ steigt, Drehmoment steigt, $n$ steigt.

Bei grösserem $\alpha$ verkleinert sich das Leitintervall, obwohl $U_a(\alpha)$ steuerbar ist.
Ab einem Punkt dominiert die Verkürzung des Leitintervalls, $I_{a,\mathrm{avr}}$ sinkt, Drehmoment sinkt,
und damit fällt die Drehzahl wieder.

Zusammenhang Winkel $\beta$:
\[
\sin\beta=\frac{E}{U_{2m}}
\quad\Rightarrow\quad
\boxed{\beta=\arcsin\!\left(\frac{E}{U_{2m}}\right)}.
\]
Mit
\[
E=k\Phi\omega=k_1 n
\]
ist $\beta$ drehzahlabhängig, was die nichtlineare Rückkopplung erklärt.

% ---------------------------------------------------------

\subsection*{c) Grenzbedingung für den Stromfluss}

\paragraph{Gesucht}
Grenze zwischen Stromfluss und kein Stromfluss.

\paragraph{Lösungsweg}
Grenzfall:
\[
\boxed{\alpha=\beta}
\]
mit
\[
\boxed{\beta=\arcsin\!\left(\frac{E}{U_{2m}}\right)}.
\]
Für $\alpha\le\beta$ gilt näherungsweise:
\[
I_a \approx 0,\quad M\approx 0.
\]

